\documentclass[11pt,a4paper]{article}
\usepackage[margin=2.2cm,headheight=16pt]{geometry}
\usepackage[spanish,es-nodecimaldot]{babel}
\usepackage{microtype}
\usepackage{xcolor}
\usepackage{booktabs}
\usepackage{longtable}
\usepackage{tabularx}
\usepackage{array}
\usepackage{enumitem}
\usepackage{fancyhdr}
\usepackage{lastpage}
\usepackage{titlesec}
\usepackage{amsmath}
\usepackage{tikz}
\usetikzlibrary{arrows.meta,positioning}
\usepackage[most]{tcolorbox}
\usepackage{hyperref}

\definecolor{MemBlue}{HTML}{1F49C4}
\definecolor{MemDark}{HTML}{161A21}
\definecolor{MemGray}{HTML}{4B5563}
\definecolor{MemLight}{HTML}{E7ECFA}
\definecolor{MemGreen}{HTML}{0D7355}
\definecolor{MemOrange}{HTML}{A06B06}
\definecolor{MemRed}{HTML}{AE382C}

\hypersetup{colorlinks=true,linkcolor=MemBlue,urlcolor=MemBlue,
  pdftitle={Informe final del proyecto MemLab 16 MiB},pdfauthor={Informe del proyecto MemLab}}
\pagestyle{fancy}
\fancyhf{}
\lhead{\sffamily\small MemLab 16 MiB}
\rhead{\sffamily\small Informe final}
\cfoot{\sffamily\small Página \thepage\ de \pageref{LastPage}}
\titleformat{\section}{\Large\bfseries\sffamily\color{MemBlue}}{\thesection}{0.7em}{}
\titleformat{\subsection}{\large\bfseries\sffamily\color{MemDark}}{\thesubsection}{0.7em}{}
\setlist{nosep,leftmargin=*}
\setlength{\parindent}{0pt}
\setlength{\parskip}{5pt}
\renewcommand{\arraystretch}{1.2}

\newtcolorbox{ejecutivo}{colback=MemLight!55,colframe=MemBlue!70!black,
  boxrule=0.7pt,arc=2mm,title=Resumen ejecutivo,fonttitle=\bfseries\sffamily}
\newtcolorbox{conclusionbox}{colback=MemGreen!6,colframe=MemGreen!70!black,
  boxrule=0.7pt,arc=2mm,title=Conclusión,fonttitle=\bfseries\sffamily}
\newcommand{\archivo}[1]{\texttt{#1}}

\begin{document}

\begin{titlepage}
  \centering
  \vspace*{2.0cm}
  {\sffamily\bfseries\fontsize{34}{39}\selectfont\color{MemBlue} MemLab 16 MiB\par}
  \vspace{0.5cm}
  {\sffamily\fontsize{23}{28}\selectfont Informe final del proyecto\par}
  \vspace{0.8cm}
  {\Large Simulador educativo de gestión de memoria principal\par}
  \vfill
  \begin{tabular}{rl}
    \textbf{Fecha de evaluación:} & 21 de septiembre de 2026\\
    \textbf{Commit base:} & \texttt{20cec7a}\\
    \textbf{Estado evaluado:} & copia de trabajo local actual\\
    \textbf{Tecnologías:} & HTML5, CSS y JavaScript\\
  \end{tabular}
  \vfill
  {\small Informe de cierre elaborado mediante inspección del código y validaciones no invasivas.\par}
\end{titlepage}

\tableofcontents
\newpage

\section{Resumen ejecutivo}

\begin{ejecutivo}
MemLab 16 MiB cumple su propósito central como herramienta didáctica para explorar
asignación contigua de memoria. La aplicación implementa tres esquemas de
particionado, tres algoritmos de ajuste, liberación y coalescencia de huecos,
compactación manual/automática, métricas y trazabilidad mediante bitácora. Su
arquitectura estática reduce las barreras de ejecución y despliegue. La evaluación
detectó una exposición de inserción HTML en la bitácora, algunos casos límite de
mensajería/configuración y ausencia de pruebas automatizadas; son oportunidades de
mejora, no impedimentos para el uso educativo local controlado.
\end{ejecutivo}

\section{Contexto y problema abordado}

Los conceptos de fragmentación y selección de particiones suelen presentarse como
diagramas estáticos. El proyecto transforma esos conceptos en un laboratorio
interactivo: el estudiante modifica límites del sistema operativo, particiones,
políticas y trabajos, y observa cómo cambian direcciones, utilización, desperdicio
y admisiones.

El espacio simulado contiene 16 MiB, equivalentes a \(2^{24}\) bytes. La porción
inicial pertenece al sistema operativo y el resto se administra mediante:

\begin{enumerate}
  \item particiones estáticas iguales;
  \item particiones estáticas de tamaños variables;
  \item particiones dinámicas con fusión y compactación.
\end{enumerate}

\section{Objetivos del proyecto}

\subsection{Objetivo general}

Proporcionar una aplicación web comprensible y manipulable que permita visualizar
la asignación contigua de memoria y contrastar sus algoritmos y formas de
fragmentación.

\subsection{Objetivos específicos observados}

\begin{itemize}
  \item representar direcciones físicas y tamaños en unidades binarias;
  \item modelar procesos con segmentos de código, datos, pila y heap;
  \item aplicar primer, mejor y peor ajuste donde corresponda;
  \item mostrar fragmentación interna y externa mediante mapa y KPI;
  \item permitir liberación, fusión de huecos y compactación;
  \item explicar cada decisión relevante en una bitácora;
  \item ejecutarse sin backend ni construcción previa;
  \item adaptarse a tema del sistema y pantallas estrechas.
\end{itemize}

\section{Alcance realizado}

\begin{longtable}{p{0.32\textwidth}p{0.18\textwidth}p{0.38\textwidth}}
\toprule
\textbf{Capacidad} & \textbf{Estado} & \textbf{Implementación observada}\\
\midrule
\endhead
Espacio de 16 MiB & Implementado & Constantes binarias y direcciones hexadecimales de 24 bits.\\
Reserva de S.O. & Implementado & Opciones de 1, 2 y 4 MiB al inicio del mapa.\\
Particiones fijas & Implementado & 2 a 16 particiones según opciones, redondeo a 4 KiB.\\
Particiones variables & Implementado & Tabla editable en KiB y zona no particionada.\\
Particiones dinámicas & Implementado & División exacta de huecos y cero fragmentación interna.\\
Tres algoritmos & Implementado & Primer, mejor y peor ajuste por recorrido lineal.\\
Liberación y fusión & Implementado & Coalescencia de huecos dinámicos contiguos.\\
Compactación & Implementado & Manual y automática al fallar por discontinuidad.\\
Métricas & Implementado & Utilización, dos fragmentaciones, mayor bloque y procesos.\\
Programas personalizados & Implementado & Nombre y cuatro segmentos en KiB.\\
Escenario pedagógico & Implementado & Carga y liberación alternada para producir huecos.\\
Persistencia & Fuera del alcance actual & Todo el estado se pierde al recargar.\\
Backend / cuentas & Fuera del alcance actual & La aplicación es completamente estática.\\
Suite de pruebas & No incluida & No hay infraestructura ni casos automatizados en el repositorio.\\
\bottomrule
\end{longtable}

\section{Solución construida}

\subsection{Estructura}

\begin{center}
\begin{tikzpicture}[
  node distance=12mm,
  box/.style={draw=MemBlue!65,rounded corners,fill=MemLight!55,
    minimum width=46mm,minimum height=12mm,align=center,font=\sffamily\small},
  arr/.style={-{Latex[length=2.3mm]},thick,color=MemBlue!75}
]
\node[box] (html) {\archivo{index.html}\\estructura y controles};
\node[box,below left=of html] (css) {\archivo{css/style.css}\\presentación y respuesta};
\node[box,below right=of html] (js) {\archivo{js/script.js}\\modelo, algoritmos y DOM};
\node[box,below=18mm of html] (browser) {Navegador\\ejecución completa};
\draw[arr] (html)--(css);
\draw[arr] (html)--(js);
\draw[arr] (css)--(browser);
\draw[arr] (js)--(browser);
\end{tikzpicture}
\end{center}

El README completa el conjunto con instrucciones y descripción conceptual. No hay
archivos de configuración, dependencias de aplicación ni compilación.

\subsection{Arquitectura funcional}

El JavaScript está encapsulado en una IIFE con estado mutable privado. Los eventos
DOM llaman reglas de negocio sincrónicas; cada cambio termina en un renderizado
integral o parcial. La interfaz se deriva del estado, por lo que el mapa, los KPI,
la tabla y la bitácora se mantienen coordinados.

La decisión de regenerar \texttt{innerHTML} es suficiente para el volumen actual y
evita complejidad de framework. A cambio, el dominio queda acoplado al DOM y es más
difícil probarlo de forma aislada.

\subsection{Modelo de memoria}

El arreglo \texttt{mem} mantiene bloques ordenados con base, tamaño, tipo y PID.
Los procesos incluyen cuatro segmentos, tamaño agregado, estado y ubicación. La
configuración activa define método, algoritmo, reserva del S.O., particiones y
compactación automática.

La configuración inicial reserva 2 MiB y crea siete particiones de 2 MiB. De los
ocho trabajos, cuatro se cargan durante el arranque. El catálogo total requiere
18\,880 KiB y la memoria de usuario predeterminada ofrece 14\,336 KiB, lo que
garantiza decisiones visibles de admisión y rechazo.

\section{Algoritmos implementados}

\subsection{Asignación}

El proceso general filtra bloques libres con capacidad suficiente y selecciona uno:

\begin{itemize}
  \item primer ajuste: primer candidato en orden de dirección;
  \item mejor ajuste: candidato de menor capacidad válida;
  \item peor ajuste: candidato de mayor capacidad válida.
\end{itemize}

En particiones fijas iguales se fuerza la primera disponible. En dinámica, el hueco
se divide para que el bloque ocupado tenga exactamente el tamaño del proceso.

\subsection{Liberación y compactación}

Una liberación estática limpia el PID, sin alterar límites. Una liberación dinámica
convierte el bloque en hueco y combina vecinos. La compactación ordena los procesos
por dirección y los desplaza hacia el final del S.O.; todo el espacio disponible
queda, si existe, en un hueco final.

La compactación automática solo se intenta cuando el total libre es suficiente
pero ningún hueco individual alcanza. Este criterio distingue correctamente falta
de capacidad total de falta de continuidad.

\subsection{Indicadores}

La utilización se calcula sobre memoria de usuario y solo suma tamaños útiles de
procesos. La fragmentación interna suma sobrantes dentro de bloques ocupados. En
dinámica, la fragmentación externa se define como libre total menos el hueco más
grande. En estáticas, el proyecto contabiliza como externa toda partición libre no
combinable, una convención pedagógica documentada en la interfaz.

\section{Experiencia de usuario}

La página organiza la experiencia en cinco niveles:

\begin{enumerate}
  \item método y acciones globales;
  \item configuración y algoritmo;
  \item KPI de lectura inmediata;
  \item mapa, detalle y tabla de trabajos;
  \item bitácora explicativa.
\end{enumerate}

El mapa usa altura proporcional al tamaño real y adapta la cantidad de texto cuando
un bloque es pequeño. Los colores diferencian procesos y segmentos; los patrones
identifican fragmentación y espacio reservado. En pantallas de hasta 900 px, la
disposición pasa a una columna.

La interfaz incorpora etiquetas ARIA, estados \texttt{aria-pressed}, foco visible,
región activa para la bitácora, tema oscuro y consideración de movimiento reducido.

\section{Metodología de evaluación}

La revisión final fue deliberadamente no invasiva:

\begin{enumerate}
  \item inventario completo del repositorio y lectura de sus cuatro archivos;
  \item trazado de las 33 funciones declaradas y sus efectos sobre el estado;
  \item revisión de fórmulas, rangos, estados y eventos;
  \item comprobación sintáctica con Node.js;
  \item servicio temporal por HTTP y solicitud de HTML, CSS y JS;
  \item comparación automática de selectores \texttt{\#id} frente al HTML;
  \item cálculo independiente de tamaños de catálogo y configuración predeterminada;
  \item registro de huellas SHA-256 de los archivos originales para verificar que
  la generación documental no los modificara.
\end{enumerate}

Se intentó preparar una validación con navegador automatizado conforme a la guía de
pruebas disponible, pero el entorno no contiene la dependencia Python Playwright.
No se instaló en el proyecto para respetar el alcance de ``solo entregables''.

\section{Resultados de validación}

\begin{longtable}{p{0.32\textwidth}p{0.17\textwidth}p{0.39\textwidth}}
\toprule
\textbf{Comprobación} & \textbf{Resultado} & \textbf{Detalle}\\
\midrule
\endhead
Sintaxis JavaScript & Aprobada & Node.js no reportó errores de análisis.\\
Disponibilidad HTTP & Aprobada & Los tres recursos ejecutables respondieron 200.\\
Tipos MIME & Aprobada & HTML, CSS y JavaScript fueron servidos con tipos correspondientes.\\
Integridad de IDs & Aprobada & 33 IDs únicos y ningún selector auxiliar faltante.\\
Catálogo de programas & Verificado & Totales recalculados a partir de los cuatro segmentos.\\
Tamaño fijo inicial & Verificado & 14 MiB / 7 = 2 MiB por partición.\\
Tabla variable inicial & Verificada & Los seis tamaños suman 14\,336 KiB.\\
Prueba end-to-end real & Pendiente & No disponible Playwright en el entorno.\\
\bottomrule
\end{longtable}

\section{Evaluación de calidad}

\begin{longtable}{p{0.24\textwidth}p{0.17\textwidth}p{0.49\textwidth}}
\toprule
\textbf{Atributo} & \textbf{Valoración} & \textbf{Fundamento}\\
\midrule
\endhead
Adecuación funcional & Alta & Cubre métodos, ajustes, métricas, escenarios y explicación de decisiones.\\
Comprensibilidad & Alta & Código comentado, nombres descriptivos y bitácora pedagógica.\\
Portabilidad & Alta & Archivos estáticos sin construcción ni backend.\\
Rendimiento & Adecuado & Recorridos lineales y render total sobre conjuntos pequeños.\\
Mantenibilidad & Media-alta & Separación por funciones, aunque modelo y DOM comparten un único módulo.\\
Accesibilidad & Media-alta & Buenas bases ARIA y CSS; falta verificación formal con herramientas/lectores.\\
Seguridad & Media & Superficie pequeña, pero existe inserción HTML en la bitácora.\\
Testabilidad & Media-baja & No hay pruebas y el estado privado está acoplado al navegador.\\
Observabilidad & Alta para docencia & La bitácora explica causas y efectos hasta 80 eventos.\\
\bottomrule
\end{longtable}

\section{Hallazgos, riesgos y limitaciones}

\subsection{Prioridad alta}

\textbf{Riesgo de XSS DOM en nombres personalizados.} El nombre se escapa al
mostrarlo en mapa, detalle y tabla, pero se concatena sin escape a un mensaje de la
bitácora, que luego se inserta mediante \texttt{innerHTML}. El riesgo es limitado
en uso local individual, pero debe corregirse antes de aceptar datos no confiables
o publicar la aplicación en un entorno compartido.

\subsection{Prioridad media}

\begin{itemize}
  \item al aumentar el tamaño del S.O., la tabla variable puede sumar más que la
  memoria disponible; el constructor omite desde la primera partición que no cabe,
  pero el texto de suma sigue considerando toda la configuración;
  \item algunos mensajes de rechazo del modo fijo pueden atribuir el fallo a
  ocupación total cuando existen particiones libres pero demasiado pequeñas;
  \item compactar una memoria totalmente llena puede registrar ``1 hueco de 0 KiB''
  aunque no se agregue un bloque de hueco;
  \item no existe persistencia ni exportación de escenario/bitácora;
  \item no hay pruebas automatizadas de regresión.
\end{itemize}

\subsection{Prioridad baja o pedagógica}

\begin{itemize}
  \item la denominación de fragmentación externa en esquemas estáticos es una
  simplificación didáctica que debe explicarse al compararla con otras fuentes;
  \item \texttt{color-mix} puede degradar la apariencia en navegadores antiguos;
  \item la lista de eventos se limita a 80 y no representa tiempo real;
  \item el sistema simula asignación contigua y no cubre paginación, segmentación
  no contigua, memoria virtual, swapping ni traducción mediante MMU.
\end{itemize}

\section{Entregables documentales}

Como parte del cierre se produjeron, sin modificar el código fuente:

\begin{longtable}{p{0.32\textwidth}p{0.55\textwidth}}
\toprule
\textbf{Entregable} & \textbf{Contenido}\\
\midrule
\endhead
Manual técnico (LaTeX y PDF) & Arquitectura, datos, algoritmos, fórmulas, catálogo de funciones, eventos, ejecución, despliegue, validación, riesgos y mantenimiento.\\
Manual de usuario (LaTeX y PDF) & Instalación, recorrido, operación paso a paso, ejercicios, solución de problemas y glosario.\\
Informe final (LaTeX y PDF) & Objetivos, alcance, solución, evaluación, resultados, calidad, riesgos y recomendaciones.\\
\bottomrule
\end{longtable}

Los tres fuentes LaTeX son autónomos y reproducibles con un motor XeLaTeX moderno
o Tectonic, sujeto a disponibilidad de los paquetes estándar declarados.

\section{Recomendaciones}

\subsection{Corto plazo}

\begin{enumerate}
  \item eliminar la inserción HTML no escapada; preferir \texttt{textContent} o
  escapar el mensaje completo de bitácora;
  \item añadir pruebas unitarias para tamaños, selección, fusión y métricas;
  \item cubrir los casos límite de mensajes y compactación sin espacio libre;
  \item validar/ajustar la tabla variable al cambiar la reserva del S.O.;
  \item documentar expresamente la convención de fragmentación externa estática.
\end{enumerate}

\subsection{Mediano plazo}

\begin{enumerate}
  \item extraer el dominio a un módulo JavaScript puro y dejar el DOM como adaptador;
  \item incorporar una suite end-to-end con escenarios reproducibles;
  \item permitir exportar/importar configuración, trabajos y bitácora;
  \item incluir un modo de comparación lado a lado entre algoritmos;
  \item ejecutar auditorías automatizadas de accesibilidad y compatibilidad móvil.
\end{enumerate}

\subsection{Extensiones académicas}

\begin{itemize}
  \item siguiente ajuste (\emph{next fit}) y listas libres ordenadas;
  \item paginación, tablas de páginas y fallos de página;
  \item segmentación con protección y compartición;
  \item costo temporal de compactación y reubicación;
  \item generación de cargas aleatorias con métricas comparativas.
\end{itemize}

\section{Integridad del código fuente}

Antes de producir los documentos se registraron las siguientes huellas:

\begin{longtable}{p{0.21\textwidth}p{0.68\textwidth}}
\toprule
\textbf{Archivo} & \textbf{SHA-256 inicial}\\
\midrule
\endhead
\archivo{README.md} & \texttt{b00fcb9c6884db5620871af1cc4c5b36}\allowbreak\texttt{fcec4d064e7b4008b878589c7f9e131a}\\
\archivo{index.html} & \texttt{67277c06878c8ee2f17b20d6b2a92f5c}\allowbreak\texttt{a0de73367ba3dcb3cf4e67f600e6b37c}\\
\archivo{css/style.css} & \texttt{083eaae30a30cc87688f7eb674c7a4732}\allowbreak\texttt{613a009e279b7b18eb6c2097c7c79f3}\\
\archivo{js/script.js} & \texttt{70cead1c39347565160049b735fcca059}\allowbreak\texttt{16dbcf603c264c4ccb0e9b61f88f3ba}\\
\bottomrule
\end{longtable}

Estas huellas sirven para comparar el estado posterior y demostrar que la tarea de
documentación creó únicamente los entregables solicitados.

\section{Conclusión final}

\begin{conclusionbox}
El proyecto está completo y coherente como simulador educativo estático de
asignación contigua. Su principal fortaleza es convertir cada operación en una
representación visual y una explicación rastreable. La base técnica es pequeña,
portable y suficientemente clara para enseñanza. Antes de una publicación abierta
o una ampliación funcional conviene corregir el manejo HTML de la bitácora,
formalizar pruebas y separar el dominio de la capa de presentación. Los manuales y
este informe dejan documentado tanto el funcionamiento actual como la ruta de
evolución recomendada.
\end{conclusionbox}

\end{document}
